%Tex root = ProbabilitàEStatistica.tex

\chapter{Introduzione}
La \textcolor{viola}{\textbf{statistica}} si occupa della \hl{raccolta dati, nella loro descrizione e analisi}, guidandoci nel trarre le conclusioni.

\section{Raccolta dati e statistica descrittiva}
La statistica permette di descrivere, sintettizzare e analizzare dati presi da un campione di dati definito.\\[0.2ex]
Altre volte, il lavoro dello statistico inizia prima che i dati siano stati ottenuti e il suo primo obiettivo consiste nell'ideare un procedimento ottimale per la loro raccolta.

\esempio{indagine statistica}{
    Immaginiamo che un docente voglia determinare quale sia il più efficace tra due diversi metodi per insegnare la programmazione a dei neofiti.

    \smallskip
    Un possibile approcio potrebbe essere:
    \begin{enumerate}
    \item dividere gli studenti in due gruppi
    \smallskip
    \item usare un metodo didattico per ciascun gruppo
    \end{enumerate}

    \smallskip
    Alla fine del corso gli studenti vengono esaminati e i punteggi dei membri dei due gruppo vengono confrontati.
    
    Se i risultati di uno dei due gruppi risultassarso notelvolemente più alti, sarebbe ragionevole pensare che il corrispondente metodo di insegnamento sia migliore.

    \smallskip
    È importante notare che per poter trarre delle conclusioni valide dei dati è essenziale che gli studenti siano divisi in modo che in nessuno dei due gruppi si vengano a trovare elementi con una maggiore predispozione alla programmazione.

    Il metodo accettato per superare questo problema consiste nel dividere "a caso" gli studenti in due gruppi.
}

A esperimento concluso, i dati devono essere raccolti e commentati.

\bigskip
La parte della \hl{statistica che si occupa di illustrare e sintettizzare i dati} è chiamata \textcolor{viola}{\textbf{statistica descrittiva}}.

\section{Inferenza statistica e modelli probabilistici}
Considerando l'esempio fatto in precedenza, dopo che la prova si è conclusa e i dati sono stati illustrati e sintetizzati, vorremmo poter trarre una conclusione su quale dei due metodi di insegnamento sia superiore.

\medskip
La parte della statistica che si occupa di questo aspetto è detta \textcolor{viola}{\textbf{inferenza statistica}}.

Per dedurre enenunciati formalmemte validi dai dati raccolti, è necessario prendere in considerazione l'influenza del caso.

\esempio{indagine statistica}{
    Potrebbe accadere che il punteggio medio dei membri del primo gruppo sia superiore, ma di poco, a quello del secondo gruppo.

    \smallskip
    \textcolor{viola}{\textbf{È corretto concludere che questa differenza sia dovuta al metodo didattico utilizzato?}}

    \smallskip
    \textcolor{viola}{\textbf{È possible che sia una casualità?}}
}

Per poter giungere a conclusioni pienamente giustificate è necessario fare alcune assunzioni sulla probabilità che i dati che vengono misurati assumano diversi valori possibili.\\[0.2ex]
L'\hl{insieme di queste ipotesi} è detto \textbf{modello probabilistico} per i dati.

A volte la natura dei dati suggerisce la forma del modello probabilistico da adottare.

\esempio{indagine statistica}{
    Consideriamo un ingegnere della produzione che vuole scoprire la frazione di circuiti integrati difettosi riscontrata con un nuovo metodo produttivo.

    \smallskip
    Egli potrebbe selezionare un certo numero di chip e testarli: il numero di quelli difettosi costituirà il dato sperimentale.

    \smallskip
    Se la scelta del campione da testare è stata fatta "a caso" è ragionevole supporre che ciascuno dei chip sarà difettoso con probabilità $\mathrm{p}$, pari alla frazione incognita di chip difettosi sul totale di quelli prodotti.
}

In altre situazioni potrebbe non essere evidente quale sia il modello di probabilità per un certo campione di dati.\\[0.2ex]
Molto spesso, un'attenta descrizione e presentazione dei dati ci permette di inferire quale sia un modello accettabile, che può eventualmente essere messo alla prova raccogliendo nuovi dati.

\section{Popolazioni e campioni}
La statistica è normalmente interessata a ottenere informazioni su un \hl{insieme completo di oggetti} che viene detto \textbf{popolazione}.\\[0.2ex]
Esso, però, è spesso troppo grande perchè sia possibile un esame esaustivo.\\[0.2ex]
In questi casi, si cerca di imparare qualcosa sulle popolazioni scegliendo e poi esaminnado dei sottogruppi di loro elementi.\\[0.2ex]
Un \hl{sottogruppo della popolazione} è detto \textbf{campione}.\\[0.2ex]
Siccome il campione deve contenere informazioni sulla popolazione complessiva, deve essere rapresentativo di quella popolazione.

\esempio{indagine statistica}{
    Consideriamo di effettuare una ricerca sulla distribuzione delle età degli abitati di un certo comune e, intervistati i primi 100 che entrano in una biblioteca, venisse trovata che l'età media è di 46.2 anni.

    \smallskip
    \textcolor{viola}{\textbf{Possiamo concludere che questa sia l'età media dell'intera popolazione?}}

    \smallskip
    Probabilmente no, perchè il campione scelto non è rappresentativo della popolazione in esame, essendo gli utenti della biblioteca più facilmente studenti e anziani che non persone in età lavorativa.
}